\chapter{Geometric pushforwards and integer multiplicities}\label{sec:geometric-integer-pushforward}
We connect the analytic construction for normal currents with the tangential formula of \cite{Sim18}, Chapter 6, Lemma 3.12 and Remark 3.13, printed pp.~174--175. The approximation argument requires weak continuity of derivative minors, rather than pointwise convergence of ambient derivatives on a surface. We supply that step first. The local area formula, the Whitney-based $C^1$ decomposition and the compact normal construction remain identified inputs; no integral-current compactness theorem is used.

\begin{lemma}[Distributional exterior products of Lipschitz differentials]\label{lem:lipschitz-differential-wedge-identity}
Let $V\subset\mathbb R^m$ be open, let $1\leq r\leq m$, and let $u^1,\ldots,u^r$ be locally Lipschitz scalar functions. Define the coefficients of their differential forms by their almost-everywhere derivatives. Then, in distributions,
\[
 d(u^1\,du^2\wedge\cdots\wedge du^r)
 =du^1\wedge\cdots\wedge du^r.
\]
For $r=1$ the expression in parentheses is $u^1$. In particular every product of Lipschitz differentials is distributionally closed.
\end{lemma}
\begin{proof}
Fix a compactly supported smooth complementary-degree test form in $V$ and work in a slightly larger relatively compact neighborhood. Mollify the functions there. Their values converge uniformly, their derivatives are uniformly bounded, and the derivatives converge in every finite $L^p$ norm on the compact neighborhood. The last fact is the approximate-identity property applied to the weak derivatives, which agree almost everywhere with the classical derivatives of a Lipschitz function; this agreement follows by one-dimensional integration by parts along coordinate lines and Fubini. All products of the finitely many derivative coefficients therefore converge in $L^1$: expand a difference of products one factor at a time and use the uniform bounds. The factors with an undifferentiated function converge in the same way. For smooth mollifications the identity is the exterior product rule and $d^2=0$. Pass to the limit after testing and integration by parts. Applying $d$ once more gives closedness, since distributional mixed partials commute. No differentiability on an ambient rectifiable support is asserted here; the domain is the Euclidean coordinate space.
\end{proof}

\begin{proposition}[Weak continuity of minors under a uniform Lipschitz bound]\label{prop:uniform-lipschitz-weak-minors}
Let $u_j,u:V\subset\mathbb R^m\to\mathbb R^q$ be locally Lipschitz. Suppose $u_j\to u$ locally uniformly and their Lipschitz constants are uniformly bounded on each relatively compact ball. Every minor of $Du_j$ of a fixed order $1\leq r\leq\min\{m,q\}$ converges weakly against compactly supported $L^1$ coefficients to the corresponding minor of $Du$. Equivalently the coefficient forms
\[
 du_j^{i_1}\wedge\cdots\wedge du_j^{i_r}
 \ \rightharpoonup^*\ 
 du^{i_1}\wedge\cdots\wedge du^{i_r}
 \quad\text{in }L^\infty_{\mathrm{loc}}.
\]
Repeated target indices simply give zero forms.
\end{proposition}
\begin{proof}
For order one, distributional integration by parts and local uniform convergence give convergence against smooth tests. The uniform derivative bound extends it to every $L^1$ test by approximation with smooth functions on a common compact neighborhood.

Induct on $r$. By Lemma \ref{lem:lipschitz-differential-wedge-identity}, test the order-$r$ form by integrating by parts onto the complementary test form, leaving
$u_j^{i_1}du_j^{i_2}\wedge\cdots\wedge du_j^{i_r}$.
The difference made by replacing the undifferentiated factor $u_j^{i_1}$ with $u^{i_1}$ tends to zero by uniform convergence and the bounded lower-order minors. With that factor fixed, induction gives convergence, since its product with a smooth test coefficient is an $L^1$ coefficient. Apply the same distributional identity to the limit. This proves distributional convergence of every order-$r$ coefficient. All these coefficients have a common local bound, from the uniform derivative bound and the determinant expansion. Smooth approximation of an arbitrary $L^1$ coefficient again gives weak-star convergence. A finite partition of unity covers the support of a test meeting several balls. This argument proves convergence of minors, not strong or pointwise convergence of the derivative matrices.
\end{proof}

\begin{lemma}[Pullback integrals converge against a fixed integrable weight]\label{lem:weighted-pullback-weak-convergence}
Under Proposition \ref{prop:uniform-lipschitz-weak-minors}, let $1\leq m\leq q$, let $a\in L^1(V)$ have compact support in $V$, and let $\omega$ be a smooth $m$-form on the target. Then
\[
 \int_V a(x)\bigl\langle\omega(u_j(x)),
       \Lambda^mDu_j(x)(e_1\wedge\cdots\wedge e_m)\bigr\rangle\,dx
 \longrightarrow
 \int_V a(x)\bigl\langle\omega(u(x)),
       \Lambda^mDu(x)(e_1\wedge\cdots\wedge e_m)\bigr\rangle\,dx.
\]
The target form need only be smooth near the common compact set containing these images.
\end{lemma}
\begin{proof}
Expand the form in the finite coordinate basis $dy^{i_1}\wedge\cdots\wedge dy^{i_m}$. Each summand is its coefficient evaluated at $u_j$ times one $m$-minor of $Du_j$. Uniform convergence of the maps makes the evaluated smooth coefficients converge uniformly on the compact support of $a$. The uniformly bounded minors control the resulting error by a constant times $\|a\|_1$ times that uniform difference. After replacing the coefficient by its limit, use Proposition \ref{prop:uniform-lipschitz-weak-minors} with the fixed $L^1$ weight $a$ times that coefficient. Summing finitely many terms proves the assertion. In dimension zero the corresponding statement is ordinary uniform convergence of functions tested against a finite signed measure; minors are not needed.
\end{proof}

\begin{lemma}[Full-rank image tangents agree with the image of the differential]\label{lem:rectifiable-full-rank-image-tangent}
Let $M\subset\mathbb R^N$ be measurable and countably $m$-rectifiable with finite $\mathcal H^m$ measure, $m\geq1$, and let $f:\mathbb R^N\to\mathbb R^q$ be Lipschitz, with $q\geq m$. Its image is countably $m$-rectifiable and has finite $\mathcal H^m$ measure. At almost every $x\in M$ with $J_m(d^Mf_x)>0$,
\[
 T_{f(x)}f(M)=d^Mf_x(T_xM).
\]
Null sets in the image are understood using measurable representatives. This is the tangent assertion preceding \cite{Sim18}, equation (3.6), p.~174.
\end{lemma}
\begin{proof}
Compose the countably many Lipschitz parametrizations of $M$ with $f$. Their images cover $f(M)$ except for the image of a Hausdorff-null set, which is null by the Lipschitz measure bound. The same bound gives finite image measure and measurability follows from compact cores.

Work on a regular $C^1$ source chart $F$ and put $u=f\circ F$. By Corollary \ref{cor:vector-c1-lusin-approximation}, countably many $C^1$ maps agree with $u$ and its derivative on measurable sets covering almost every parameter point. On the sets where the derivative has rank $m$, subdivide further among neighborhoods where the corresponding $C^1$ map is an injective immersion. The inverse function theorem applied to a nonzero coordinate minor supplies these neighborhoods; their images are embedded $C^1$ manifolds. The chain rule in Theorem \ref{thm:rectifiable-tangential-differential} identifies their tangent images with $d^Mf(T_xM)$ at the retained points.

The measurable image piece lies in both this smooth image manifold and $f(M)$. Corollary \ref{cor:rectifiable-common-tangents} makes their tangents agree outside an image-null set. Enlarge the exceptional image set to a Borel Hausdorff-null set before taking its preimage. Its preimage on the positive-Jacobian part is domain-null: apply Theorem \ref{thm:rectifiable-weighted-area-local} with its indicator to get integral of $J_m(d^Mf)$ zero there, and use its strict positivity. Thus the tangent equality holds at almost every full-rank domain point. Countably many pieces remove only one null set. The rank-deficient domain part has null image by the same area formula and the indicator argument; it creates no additional almost-everywhere image tangent. This justifies both the domain and image exceptional-set conventions.
\end{proof}

\begin{theorem}[Signed image multiplicity of an integer rectifiable current]\label{thm:integer-rectifiable-geometric-image}
Let $T=\llbracket M,\theta,\tau\rrbracket$ be an integer rectifiable $m$-current of finite mass in $\mathbb R^N$, with $m\geq1$, and let $f:\mathbb R^N\to\mathbb R^q$ be $L$-Lipschitz. Define
\[
 G_fT(\omega)=\int_M\theta(x)
 \bigl\langle\omega(f(x)),\Lambda^m d^Mf_x(\tau(x))\bigr\rangle
 \,d\mathcal H^m(x).
\]
This is an integer rectifiable current, with $\mathbf M(G_fT)\leq L^m\mathbf M(T)$. If $q\geq m$, choose a measurable unit orientation $\eta$ of the approximate tangent plane of $f(M)$. At full-rank points write
\[
 \Lambda^m d^Mf_x(\tau(x))
 =\sigma(x)J_m(d^Mf_x)\eta(f(x)),\qquad \sigma(x)\in\{-1,1\}.
\]
For almost every target point, its signed multiplicity and absolute multiplicity are
\[
 b(y)=\sum_{x\in M\cap f^{-1}(y),\,J_m(d^Mf_x)>0}
           \sigma(x)\theta(x),
 \qquad N(y)=|b(y)|.
\]
These sums are absolutely finite almost everywhere. They give the induced orientation $\operatorname{sign}(b)\eta$ on $\{N>0\}$. This is the content of \cite{Sim18}, Lemma 3.12 in the finite-mass global setting; no assertion about boundary is needed here.
\end{theorem}
\begin{proof}
Since $\theta\geq1$, finite mass gives $\mathcal H^m(M)<\infty$. The integrand is bounded in absolute value by $L^m\theta\|\omega\|_{*,\infty}$, so the formula defines a finite-mass current. If $q<m$, every $m$-vector in the target vanishes and the result is zero. Otherwise Lemma \ref{lem:rectifiable-full-rank-image-tangent} gives the asserted alignment at almost every full-rank point. A Borel representative of a measurable orientation can be chosen by partitioning a Borel representative of the rectifiable image among its $C^1$ pieces and taking each chart's orientation; null exceptional points do not affect the formula.

The weighted area formula gives
\[
 \int_{\mathbb R^q}\sum_{x\in M\cap f^{-1}(y),\,J_m(d^Mf_x)>0}
                 \theta(x)\,d\mathcal H^m(y)
 =\int_M\theta J_m(d^Mf)\,d\mathcal H^m
 \leq L^m\mathbf M(T).
\]
Thus the nonnegative sum is finite almost everywhere; because its terms are positive integers, there are only finitely many contributing points there. Apply the same weighted formula to the positive and negative parts of the oriented integrand, or componentwise to the finite coordinate coefficients. It rewrites $G_fT(\omega)$ as $\int_{f(M)}b\langle\omega,\eta\rangle\,d\mathcal H^m$. The rank-deficient part has zero exterior differential and contributes zero. The integer-valued $b$ is measurable and integrable; orient where it is negative by reversing $\eta$ and discard its zero set. This is precisely an integer rectifiable representation, with mass at most the displayed unsigned sum. Overlaps are combined with their signs before taking absolute value.
\end{proof}

\begin{proposition}[The analytic and tangential pushforward formulas agree]\label{prop:normal-geometric-pushforward-agreement}
Let $T$ be compactly supported, normal and rectifiable, with a representation
\[
 T(\omega)=\int_M a(x)\langle\omega(x),\tau(x)\rangle\,d\mathcal H^m(x),
\]
where $M$ is countably $m$-rectifiable with locally finite $\mathcal H^m$ measure, $a$ is real and integrable, and $\tau$ is a measurable unit tangent orientation. For every globally Lipschitz $f$, its analytic pushforward in Theorem \ref{thm:compact-normal-lipschitz-pushforward} equals the tangential integral obtained by replacing $\tau$ by $\Lambda^m d^Mf(\tau)$ and evaluating the target form at $f(x)$. Integer multiplicities are not required for this agreement.
\end{proposition}
\begin{proof}
For $m=0$ both definitions are the pushforward of a finite signed measure. If the target dimension is smaller than $m$, both actions are zero because there are no nonzero target $m$-forms. Otherwise, for $m\geq1$, mollify $f$ to smooth maps $f_j$ with a common global Lipschitz bound. The analytic definition is their weak limit. Partition $M$ into disjoint measurable pieces in relatively compact regular $C^1$ charts $F:V\to\mathbb R^N$. Subdivide so $DF$ is bounded and its least singular value is bounded below. Choose the chart orientation and absorb its agreement sign with $\tau$ into the scalar coefficient. The area formula then rewrites the contribution of a piece as
\[
 \int_V c(z)\bigl\langle\omega(f_j(F(z))),
           \Lambda^mD(f_j\circ F)(z)(e_1\wedge\cdots\wedge e_m)
                 \bigr\rangle\,dz,
\]
where $c$ is the pullback of $a$ times that sign and the indicator of the piece. It is integrable because the chart Jacobian is bounded below; use compact exhaustion within the chart if its support approaches the chart boundary. The maps $f_j\circ F$ converge uniformly on compact subsets and have uniformly bounded local Lipschitz constants. Lemma \ref{lem:weighted-pullback-weak-convergence} gives the limit on each compact portion. The intrinsic tangential chain rule identifies the limit with the claimed geometric expression.

A finite union of compact chart portions captures all but arbitrarily small $\int_M|a|\,d\mathcal H^m$. Both the smooth and limiting integrands are bounded there by $L^m\|\omega\|_{*,\infty}|a|$, using the norm bound for the intrinsic differential. Hence the remaining tail is uniformly small. First pass to the limit on the finite union, then let its omitted mass tend to zero. This proves equality on every test form. In particular it applies to a smooth scalar cutoff of an integral current even though that cutoff need not have integer multiplicities.
\end{proof}

\begin{lemma}[Localization of a proper Lipschitz map on a normal current]\label{lem:local-proper-normal-pushforward}
Let $U\subset\mathbb R^N$ and $W\subset\mathbb R^q$ be open. Let $T$ be locally normal in $U$ and $f:U\to W$ locally Lipschitz and proper on $\operatorname{spt}T$. There is a locally normal pushforward defined on each compactly supported target test by the compact normal construction with a source cutoff. It is independent of the cutoff and local extension, has support in $f(\operatorname{spt}T)$, and satisfies
\[
 \partial(f_\#T)=f_\#(\partial T).
\]
For $m=0$ the boundary assertion has the usual zero-boundary meaning.
\end{lemma}
\begin{proof}
For a compact target set $K\Subset W$ choose a relatively compact open target neighborhood $O$ of $K$. Properness makes $C=\operatorname{spt}T\cap f^{-1}(\overline O)$ compact in $U$. Choose $\chi\in C_c^\infty(U)$ equal to one on a neighborhood of $C$. The cutoff $Q=T\llcorner\chi$ is a compact normal current in the whole ambient Euclidean space by Proposition \ref{prop:current-cutoff-boundary}, extending it by zero outside $U$.

Choose a globally Lipschitz map $F:\mathbb R^N\to\mathbb R^q$ equal to $f$ near $\operatorname{spt}\chi$. Such an extension exists: multiply $f$ by a smooth compactly supported function in $U$ equal to one on a slightly larger neighborhood and extend the product by zero. The resulting map is locally Lipschitz everywhere with a common derivative/Lipschitz bound, by a finite cover of its compact support; subdividing line segments proves a global Lipschitz bound. Its target may leave $W$, which causes no difficulty because target tests are extended by zero outside $W$. Define the value of the desired pushforward on tests supported in $K$ by $F_\#Q$.

Two extensions agree near the support of $Q$, so Proposition \ref{prop:normal-pushforward-approximation-independence} makes their pushforwards agree. For two cutoffs use one extension on the union of their supports. Their difference current $T\llcorner(\chi_1-\chi_2)$ is compact and normal; its support has no point mapping into $O$, since both cutoffs equal one near $C$. The compact support theorem therefore makes its pushforward vanish on $O$. This proves agreement on every test in $K$. It also supplies linearity by choosing common cutoffs for finitely many tests. On a fixed target compact set, the mass and boundary bounds of the compact normal construction give local normality and continuity. Tests supported off $f(\operatorname{spt}T)$ vanish by the same support argument. That image is relatively closed in $W$ by properness.

For $m\geq1$, finally use
\[
 \partial Q=(\partial T)\llcorner\chi-T\llcorner d\chi.
\]
The first term is compact normal. The second term is also compact normal, being the difference of that term and $\partial Q$; its support is contained in $\operatorname{spt}T\cap\operatorname{spt}d\chi$ and therefore maps outside $O$. Its pushforward vanishes on target tests supported in $K$. Boundary commutation for $Q$ consequently gives the desired identity there. Since $\operatorname{spt}\partial T\subset\operatorname{spt}T$, the same properness applies to the boundary. In dimension zero both boundaries vanish. Exhausting target compact sets proves the statement. This proof uses normal cutoffs without claiming that those cutoffs preserve integer multiplicity.
\end{proof}

\begin{proposition}[The exact mass lost by cancellation]\label{prop:pushforward-cancellation-mass}
In Theorem \ref{thm:integer-rectifiable-geometric-image}, with $q\geq m$, put
\[
 P(y)=\sum_{\sigma(x)=1,\,f(x)=y,\,J_m(d^Mf_x)>0}\theta(x),
 \qquad
 Q(y)=\sum_{\sigma(x)=-1,\,f(x)=y,\,J_m(d^Mf_x)>0}\theta(x).
\]
Then the current image has density $|P-Q|$, while the weighted unoriented area image has density $P+Q$. Their difference is the nonnegative measure
\[
 2\min\{P,Q\}\,\mathcal H^m\llcorner f(M).
\]
They agree if and only if almost every fiber has contributions of at most one sign. This is the cancellation criterion in \cite{Sim18}, Remark 3.13; the criterion is relative to any chosen measurable image orientation and is independent of that choice.
\end{proposition}
\begin{proof}
The signed multiplicity formula gives $b=P-Q$. The weighted area formula gives the unsigned density $P+Q$. Both are finite almost everywhere and integrable by Theorem \ref{thm:integer-rectifiable-geometric-image}. For nonnegative real numbers, $P+Q-|P-Q|=2\min\{P,Q\}$. Apply this pointwise and integrate over each Borel target set. The difference measure is zero precisely when the minimum vanishes almost everywhere. Reversing the chosen target orientation exchanges $P$ and $Q$ without changing the minimum or either density. Thus the criterion records cancellation of actual oriented sheets, rather than a failure of the area formula.
\end{proof}
